%%% MAKE DOC %%%
\documentclass[11pt,oneside]{scrbook} %scrartcl
% \documentclass[11pt,oneside]{report}
% \documentclass[10pt,twocolumn,landscape]{scrartcl} 
% \setlength{\columnsep}{.5in}

%%% FRONTMATTER %%%

\title{CODING.CARE: GUIDEBOOKS FOR INTERSECTIONAL AI}
\author{Sarah Ciston}
% \majorfield{CINEMATIC ARTS (MEDIA ARTS AND PRACTICE)}
% \submitdate{May 2024}
% \subtitle{}
\setcounter{tocdepth}{1} 
\usepackage{tocloft}
\renewcommand{\cftchapdotsep}{\cftdotsep}

%%% LAYOUT %%%
\usepackage{geometry}
\geometry{letterpaper, tmargin=1in, bmargin=1in, lmargin=1in, rmargin=1in }

% \setlength{\parskip}{8pt plus 2pt minus 1pt}
\raggedright
\linespread{2}
\usepackage{setspace}
% \renewcommand*{\sectionformat}[1]{} %remove section numbering


%%% TYPESET %%%
\usepackage[utf8]{inputenc} 
\usepackage[tt=false]{libertine}
% \usepackage{fontspec}
% \usepackage{emoji}
% \setromanfont{Noto Sans}
% \setmonofont{Noto Sans Mono}
% \setemojifont{Apple Color Emoji} % must use syntax \emoji{%nameofemoji%}
\usepackage[normalem]{ulem} %for strikethrough

% \allsectionsfont{\ttfamily}
% \usepackage{koma-script}

% \usepackage{scrlayer-scrpage}
% \addtokomafont{}{\ttfamily}
% \addtokomafont{chapter}{\ttfamily}

% \usepackage[english]{babel}
% \usepackage[autostyle]{csquotes}

%%% LINKS %%%
\usepackage[breaklinks=true]{hyperref}
\hypersetup{colorlinks,%
citecolor=black,%
filecolor=black,%
linkcolor=black,%
urlcolor=black}
$if(url)$
\usepackage{url}
$endif$


%%% GRAPHICS %%%

\usepackage{graphicx}
% \setkeys{Gin}{width=0.75\linewidth}
\makeatletter
\def\maxwidth{\ifdim\Gin@nat@width>\linewidth\linewidth
\else\Gin@nat@width\fi}
\makeatother
\let\Oldincludegraphics\includegraphics
\renewcommand{\includegraphics}[1]{\Oldincludegraphics[width=\maxwidth]{#1}}

\usepackage[hybrid]{markdown}
% \usepackage{emoji}


%%% CITATION FORMATTING %%%
$if(csl-refs)$
% definitions for citeproc citations
\NewDocumentCommand\citeproctext{}{}
\NewDocumentCommand\citeproc{mm}{%
  \begingroup\def\citeproctext{#2}\cite{#1}\endgroup}
\makeatletter
 % allow citations to break across lines
 \let\@cite@ofmt\@firstofone
 % avoid brackets around text for \cite:
 \def\@biblabel#1{}
 \def\@cite#1#2{{#1\if@tempswa , #2\fi}}
\makeatother
\newlength{\cslhangindent}
\setlength{\cslhangindent}{1.5em}
\newlength{\csllabelwidth}
\setlength{\csllabelwidth}{3em}
\newenvironment{CSLReferences}[2] % #1 hanging-indent, #2 entry-spacing
 {\begin{list}{}{%
  \setlength{\itemindent}{0pt}
  \setlength{\leftmargin}{0pt}
  \setlength{\parsep}{0pt}
  % turn on hanging indent if param 1 is 1
  \ifodd #1
   \setlength{\leftmargin}{\cslhangindent}
   \setlength{\itemindent}{-1\cslhangindent}
   \setlength{\parindent}{0pt}
  \fi
  % set entry spacing
  \setlength{\itemsep}{#2\baselineskip}}}
 {\end{list}}
\usepackage{calc}
\newcommand{\CSLBlock}[1]{\hfill\break#1\hfill\break}
\newcommand{\CSLLeftMargin}[1]{\parbox[t]{\csllabelwidth}{\strut#1\strut}}
\newcommand{\CSLRightInline}[1]{\parbox[t]{\linewidth - \csllabelwidth}{\strut#1\strut}}
\newcommand{\CSLIndent}[1]{\hspace{\cslhangindent}#1}
$endif$
\setlength{\parindent}{3em}
\def\tightlist{}

% \newlength{\bibitemsep}\setlength{\bibitemsep}{.0\baselineskip plus .05\baselineskip minus .05\baselineskip}
% \newlength{\bibparskip}\setlength{\bibparskip}{0pt}
% \let\oldthebibliography\thebibliography
% \renewcommand\thebibliography[1]{%
%   \oldthebibliography{#1}%
%   \setlength{\parskip}{\bibitemsep}%
%   \setlength{\itemsep}{\bibparskip}%
% }

% hyphens, widows, orphans
\brokenpenalty 10000
\clubpenalty 10000         % Disallow club (orphan) lines
\displaywidowpenalty 10000 % Disallow widow lines before display
\widowpenalty 10000        % Disallow widow lines
\footskip   0.5 true in           % page number must be 0.5" from bottom

%%%%%%%%%%%%%%%%%%%%%%%%%%%%%%%%%%

% \newenvironment{preface}{%
%   \newpage
%   \pagenumbering{roman}%
%   \pagestyle{empty}%
%   \begin{singlespace}%
%       \maketitle
%       \setcounter{page}{1}%
%       \addtocounter{page}{\@startpageA}% Advance page counter for (signature
%                                     %   and) title.  At USC, the signature
%  \newpage                              %   page is not counted!
% %    \ifx\@volume\@empty\else % Add Volume ID to contents if present.
% %      \addtocontents{toc}{\protect\contentsline{part}{\@volume}{}}%
% %    \fi
%   \end{singlespace}%
%   \pagestyle{plain}}{%
%   \newpage
%   \pagenumbering{arabic}%
%   \pagestyle{plain}%
%   \setcounter{page}{\@startpageB}}

%%
%% Make the Dissertation titlepage
%%
% \def\@maketitle{%

%%% CONTENT %%%

\begin{document}

% \begin{preface}
% \newpage
\pagenumbering{roman}%
% PREFACE STUFF % ACKNOWLEDGMENTS
\pagestyle{empty}%
\begin{singlespace}%
    \hypersetup{hidelinks}
  \large  % title page should be 11pt or 12pt.  this should do.
  % \begin{titlepage}%
    \setcounter{page}{0}%
    \null
    % \vfill%\vfill
    \vskip 0.55 true in
    \begin{center} 
    \begin{doublespace}
        CODING.CARE: \par
    %   {\@title \par}%         Set title the same size as "FACULTY OF ..."
      % \ifx\@subtitle\@empty\else
        % \vskip 0.16 true in
        GUIDEBOOKS FOR INTERSECTIONAL AI
        % \@subtitle
      % \fi
      % changed to 0.16 by Marcos
      \vskip 0.3 true in
      by%
      \vskip 0.3 true in
      % remove tabular environment to get it aligned within
      % bounding box ---calvin 2019-04-04
      Sarah Ciston\\*[.66 true in]%
      %\rule{5.8 true in}{0.010 true in}
      % remove the horizontal rule per Manuscript formatting guidelines
      % ---calvin 2020-06-12
      \par
      \vskip 1 true in
      % Added by Marcos
    \end{doublespace}
    %\singlespace
      A Dissertation Presented to the\\%*[0.1 true in]%
      FACULTY OF THE USC GRADUATE SCHOOL\\%*[0.1 true in]%
      UNIVERSITY OF SOUTHERN CALIFORNIA\\%*[0.1 true in]%
      In Partial Fulfillment of the\\%*[0.1 true in]%
      Requirements for the Degree\\%*[0.1 true in]%
      DOCTOR OF PHILOSOPHY\\%*[0.1 true in]%
    %   \@majorfield
      CINEMATIC ARTS (MEDIA ARTS AND PRACTICE)\\%*[1.0 true cm]%
      %\vfill
      \vskip 1 true in
    %   \@submitdate
      May 2024%   should be just month and year
    \end{center}%
    \par
    \vfill
          \noindent Copyright~ 2024 \hfill Sarah Ciston
  % \end{titlepage}}
        % \maketitle
    \end{singlespace}
    \setcounter{page}{1}
    % \pagestyle{plain}%

\newpage
    % \vskip -3em
    \chapter*{Acknowledgments}
    \addcontentsline{toc}{chapter}{Acknowledgments}
% \include{path}

\newpage
\renewcommand{\contentsname}{Table of Contents}
    \begin{singlespace}
    % \vskip -3em
      \tableofcontents
    \end{singlespace}

% \newpage
    % \vskip -3em    
    % \listoffigures

% \pagestyle{plain}%
% \include{path}
  % \vskip -3em
\newpage
\chapter*{Abstract}
\addcontentsline{toc}{chapter}{Abstract}

% MOVE TO A PANDOC'D MARKDOWN > LATEX DOCUMENT I IMPORT HERE

% \import{/abstract.tex}

Critical AI researchers see the urgent need to understand and rethink how AI systems are defined, developed, regulated, and mitigated. \emph{Coding.Care: Guidebooks for Intersectional AI} argues that implementing critical approaches into AI systems more broadly requires building inviting, inclusive spaces where more people can engage creatively and critically with each other and with machine learning techniques as malleable materials. The project demonstrates craft-based, process-oriented approaches to AI that can help meet this challenge. 

The dissertation presents guides for fostering critical–creative coding communities as radical spaces of belonging—activated by an ethos and a politics modeled by queer and trans* communities that embraces radical difference. This creates a basis for deep interdisciplinary thought, interrogation of formative principles, and an openness to co-creation and alternative forms necessary to reimagine AI. 

Artists, activists, scholars, and technologists can recast their relationships to emerging technologies by reframing them as crafts like crochet—deflating AI hype, lowering barriers to learning, and emphasizing sustainability and process. Thinking craft-as-technology honors the inherited knowledge of many outsider communities. Thinking technology-as-craft provides a framework to implement those theories, ethics, and tactics as intersectional critical AI. 

\emph{Coding.Care} collects four publications that enact the strategies it theorizes. Because different modalities can better address and unite a wider range of communities, these works make interventions with different audiences by relying upon different voices, forms, formats, and media to address different stages of the sociotechnical pipelines that produce algorithmic systems like generative AI. 

"Coding.Care: Field Notes for Making Friends with Code" describes critical–creative programming approaches founded in the belief that anyone can contribute to the future of digital systems and that we all have skills to teach each other. It gives courage to pick up unfamiliar tools, find resources to kick off a new programming project, pose questions critically, or solve problems creatively. It asks: How do we code with more care? How do we encode more care into our lives? How are these connected? It supports building or joining cooperative, interdisciplinary communities for co-learning coding. This guide addresses reluctant or would-be programmers (of any age) and and potential group leaders, with a warm and friendly pocket guide, at the moment where they might intervene with critical or imaginative software creation. 

"A Critical Field Guide for Working with Machine Learning Datasets" offers practical guidance for conscientious dataset stewardship. It combines critical AI theories and technical data science concepts, explained in accessible language. It addresses journalists, students, scholars, activists, artists, and anyone starting to work with existing machine learning datasets, in the form of an instructional guidebook that combines approachable techniques with critical thinking questions, at the point when they are choosing, using, and maintaining datasets as the foundation for machine learning tasks. It is paired with the "Inclusive Datasets Research Guide," an online resource written for USC Libraries, which addresses a diverse student population who are also beginning to work with datasets. 

Both of these texts apply concepts from the "Intersectional AI Toolkit," which argues that anyone should be able to understand what AI is and help shape what AI ought to be. The Toolkit's co-authored zines are accessible guides to both AI and intersectionality. They find common vocabularies to connect diverse communities around AI's urgent questions. Its online resources learn from legacies of queer, feminist, antiracist, anticolonialist, and antiablest theories, ethics, and tactics, showing how established but marginalized tactics are necessary for reimagining more critical and ethical machine learning. The Toolkit addresses anyone who wants to understand the automated systems that impact them by using public workshops, zines, and digital resources in order to describe key concepts and processes of machine learning through critical lenses. 

As interstices among these three texts, a collection of five short lyric essays imagine dialogues with five 20th century artists, asking how the artists' analog material practices might act as pre-responses to the contemporary digital concerns raised across \emph{Coding.Care}. "Codes for (Un)Raveling," "Codes for (Un)Limiting, "Codes for (Un)Forming," "Codes for (Un)Living," and "Codes for (Un)Knowing" approach the lived experience of an algorithmic era from oblique angles. Unlike the other texts, the tone of the "Codes" essays addresses nonpractitioners on a more affective, aesthetic register, meant for reflection on the impact of sociotechnical systems as they entangle with individuals and marginalized groups. 

The guides in \emph{Coding.Care} deal with machine learning datasets, intersectional AI, coding communities, and embodied algorithmic art. Together, they want to meet readers where they are—as non-academics or those bridging into new fields, looking for a common vocabulary to engage conscientiously with the urgent concerns of AI systems. \emph{Coding.Care} offers an expansive invitation to deepen interdisciplinary conversation, apply intersectional approaches, and rework AI systems from critical–creative–caring perspectives.

\newpage
% \vskip -3em
\pagenumbering{arabic}%
\setcounter{page}{1}
% \end{preface}

\chapter{Introduction}

$body$
% \include{path}

\newpage
% \chapter*{References}
% \addcontentsline{toc}{chapter}{References}


% \addcontentsline{toc}{chapter}{\href{https://coding.care/unraveling.html}{Codes for (Un)Raveling}}

% \addcontentsline{toc}{chapter}{\href{https://coding.care/guide.html}{Coding.Care: Field Notes for Making Friends with Code}}

% \addcontentsline{toc}{chapter}{\href{https://coding.care/unlimiting.html}{Codes for (Un)Limiting}}

% \addcontentsline{toc}{chapter}{\href{https://knowingmachines.org/critical-field-guide}{A Critical Field Guide to Working with Machine Learning Datasets}}

% \addcontentsline{toc}{chapter}{\href{https://coding.care/unforming.html}{Codes for (Un)Forming}}

% \addcontentsline{toc}{chapter}{\href{https://libguides.usc.edu/inclusive-datasets}{Inclusive Datasets Research Guide}}

% \addcontentsline{toc}{chapter}{\href{https://coding.care/unliving.html}{Codes for (Un)Living}}

% \addcontentsline{toc}{chapter}{\href{https://intersectionalai.com}{Intersectional AI Toolkit}}

% \addcontentsline{toc}{chapter}{\href{https://coding.care/unknowing.html}{Codes for (Un)Knowing}}

% \addcontentsline{toc}{chapter}{References}

\chapter{\href{https://coding.care/critintro.html}{TRANS: Crafting Trans*formative Systems}}

http://coding.care/

\chapter{\href{https://coding.care/unraveling.html}{*: Codes for (Un)Raveling}}

http://coding.care/

\chapter{\href{https://coding.care/guide.html}{CRAFT: Coding.Care: Field Notes for Making Friends with Code}}

http://coding.care/

\chapter{\href{https://coding.care/unlimiting.html}{*: Codes for (Un)Limiting}}

http://coding.care/

\chapter{\href{https://knowingmachines.org/critical-field-guide}{FORMATIVE: A Critical Field Guide to Working with Machine Learning Datasets}}

http://coding.care/

\chapter{\href{https://coding.care/unforming.html}{*: Codes for (Un)Forming}}

http://coding.care/

\chapter{\href{https://libguides.usc.edu/inclusive-datasets}{FORMATIVE: Inclusive Datasets Research Guide}}

http://coding.care/

\chapter{\href{https://coding.care/unliving.html}{*: Codes for (Un)Living}}

http://coding.care/

\chapter{\href{https://intersectionalai.com}{SYSTEMS: Intersectional AI Toolkit}}

http://coding.care/

\chapter{\href{https://coding.care/unknowing.html}{*: Codes for (Un)Knowing}}

http://coding.care/

\end{document}